\chapter{拉格朗日量、场与局域粒子过程}

玻色子场、费米子场和电磁场都服从波动方程。这些方程涉及场本身以及它们对空间和时间的导数。人们可以为每个场分别列出一个运动方程，但拉格朗日量将整个系统汇集为一个紧凑的表达式。将欧拉--拉格朗日规则应用于该表达式，就能一次性得到所有场的方程。

这一构造与稳定作用量密切相关。这里直接的任务更为具体：确定场的拉格朗日量依赖于什么，用它得到几个经典方程，然后以量子力学的方式重新解释同样的场表达式。在量子理论中，拉格朗日量支配粒子的运动和相互作用，并最终决定碰撞的振幅和截面。

\section{场方程的紧凑编码}

对于一个标量场，拉格朗日密度依赖于该场及其一阶导数，
\begin{equation}
\mathcal L
=
\mathcal L\!\left(\phi,\partial_t\phi,\partial_x\phi,\ldots\right)
=
\mathcal L\!\left(\phi,\partial_\mu\phi\right).
\end{equation}
在计算时，明确列出时间导数和空间导数往往更有用；记号\(\partial_\mu\phi\)以相对论形式将同一组导数打包起来。

\begin{figure}[htbp]
\centering
\includegraphics[width=0.82\textwidth]{lecture_09_figure_03.png}
\caption{标量场拉格朗日量对场及其时空导数的依赖关系。}
\lectureframe{00:03:56}
\end{figure}

场的欧拉--拉格朗日方程为
\begin{equation}
\frac{\partial}{\partial t}
\frac{\partial\mathcal L}{\partial\dot\phi}
+
\frac{\partial}{\partial x}
\frac{\partial\mathcal L}{\partial\phi_x}
+
\frac{\partial}{\partial y}
\frac{\partial\mathcal L}{\partial\phi_y}
+
\frac{\partial}{\partial z}
\frac{\partial\mathcal L}{\partial\phi_z}
=
\frac{\partial\mathcal L}{\partial\phi},
\end{equation}
其中
\begin{equation}
\dot\phi=\frac{\partial\phi}{\partial t},
\qquad
\phi_x=\frac{\partial\phi}{\partial x},
\qquad
\phi_y=\frac{\partial\phi}{\partial y},
\qquad
\phi_z=\frac{\partial\phi}{\partial z}.
\end{equation}
如果\(\mathcal L\)包含多个场，则对每个场各应用一次该规则。结果就是运动方程组。

最简单的标量例子是
\begin{equation}
\begin{split}
\mathcal L
&=
\frac12\dot\phi^2
-\frac12\phi_x^2
-\frac12\phi_y^2
-\frac12\phi_z^2
-V(\phi)
\\
&=
\frac12\dot\phi^2-\frac12(\boldsymbol\nabla\phi)^2-V(\phi).
\end{split}
\end{equation}
空间导数项与时间导数项符号相反。逐项求导得到
\begin{align}
\frac{\partial\mathcal L}{\partial\dot\phi}&=\dot\phi,
&
\frac{\partial\mathcal L}{\partial\phi_i}&=-\phi_i,
&
\frac{\partial\mathcal L}{\partial\phi}&=-V'(\phi),
\end{align}
因此
\begin{equation}
\ddot\phi-\nabla^2\phi=-V'(\phi).
\end{equation}

\begin{figure}[htbp]
\centering
\includegraphics[width=0.74\textwidth]{lecture_09_figure_04.png}
\caption{欧拉--拉格朗日方程、标量场拉格朗日量，以及由\(V(\phi)\)生成的场方程。}
\lectureframe{00:08:23}
\end{figure}

这就是由标量拉格朗日量生成的方程。\editorialnote{所保留的画面在左边缘被裁切；所显示的补全内容由紧邻的口头推导确定。}

这里\(V(\phi)\)是场的势能密度。\(V\)中的常数项对该方程没有影响，因为其导数为零。线性势在代数上是可能的，但它没有下界：场的能量可以趋向负无穷。因此第一个稳定的例子是
\begin{equation}
V(\phi)=\frac12m^2\phi^2,
\qquad
V'(\phi)=m^2\phi,
\end{equation}
它给出
\begin{equation}
\ddot\phi-\nabla^2\phi=-m^2\phi.
\end{equation}

为了揭示\(m\)的含义，代入一个平面波，
\begin{equation}
\phi(\mathbf x,t)=\phi_0e^{i\omega t-i\mathbf k\cdot\mathbf x}.
\end{equation}
于是
\begin{equation}
\ddot\phi=-\omega^2\phi,
\qquad
\nabla^2\phi=-\mathbf k^2\phi,
\end{equation}
因此波动方程变为
\begin{equation}
-\omega^2\phi+\mathbf k^2\phi=-m^2\phi.
\end{equation}
消去\(\phi\)后，
\begin{equation}
\omega^2=\mathbf k^2+m^2.
\end{equation}
在\(\hbar=c=1\)的单位制下，频率就是能量，波数就是动量。因此
\begin{equation}
E^2=\mathbf p^2+m^2,
\end{equation}
或者，恢复\(c\)的幂次后，\(E^2=\mathbf p^2c^2+m^2c^4\)。参数\(m\)是该场一个量子的质量。

相对论不变性可以在紧凑表达式这一层面上施加。如果\(\mathcal L\)是洛伦兹标量，那么它的运动方程就是洛伦兹不变的。时间导数与空间导数之间的相对符号正是狭义相对论中熟悉的时空符号。

\section{非线性与耦合经典场}

场的更高次幂引入了一种新的行为。如果
\begin{equation}
V(\phi)=\frac12m^2\phi^2+g\phi^3,
\end{equation}
那么
\begin{equation}
\ddot\phi-\nabla^2\phi=-m^2\phi-3g\phi^2.
\end{equation}
拉格朗日量中二次的项给出场的线性方程。三次及更高次的项给出非线性方程。

这一区别在波包的运动中清晰可见。对于线性方程，一个向右运动的波包和一个向左运动的波包可以相加，构成另一个解。它们重叠并干涉一段时间，然后彼此穿过，并以原来的形状出现。有了非线性项，波包会相互影响。它们可以散射、变形，或分裂成更复杂的波。线性波的暂时干涉不是相互作用；波包发生持久的变化才标志着非线性。

同样的模式出现在两个标量场的情形中。令
\begin{align}
\mathcal L
&=
\frac12(\partial_\mu\phi)(\partial^\mu\phi)
+\frac12(\partial_\mu\rho)(\partial^\mu\rho)
-V(\phi,\rho),
\\
V(\phi,\rho)
&=
\frac12m^2\phi^2+\frac12M^2\rho^2+\rho\phi^2.
\end{align}
没有最后一项时，两个方程是独立的，
\begin{equation}
\ddot\phi-\nabla^2\phi=-m^2\phi,
\qquad
\ddot\rho-\nabla^2\rho=-M^2\rho.
\end{equation}
它们描述两种粒子的量子，质量分别为\(m\)和\(M\)。包含混合项后，
\begin{align}
\ddot\phi-\nabla^2\phi&=-m^2\phi-2\rho\phi,
\\
\ddot\rho-\nabla^2\rho&=-M^2\rho-\phi^2.
\end{align}
每个方程现在都包含另一个场。这些方程是耦合且非线性的：一个\(\phi\)波包会散射一个\(\rho\)波包，反之亦然。导数项支配能量--动量关系，二次的无导数项设定质量，混合的高阶项则编码相互作用。这就是基本的经典模式。

\section{狄拉克场与标量耦合}

对于狄拉克场，方程在时间上是一阶的。在一维简写中，并理解到可推广至三个空间方向，其项结构为
\begin{equation}
i\dot\Psi
\simeq
i\alpha\frac{\partial\Psi}{\partial x}+m\beta\Psi,
\qquad
i\boldsymbol\alpha\cdot\boldsymbol\nabla\Psi
=i\sum_i\alpha_i\frac{\partial\Psi}{\partial x_i}.
\end{equation}
矩阵\(\alpha_i\)和\(\beta\)作用于\(\Psi\)的分量。相应的拉格朗日量包含一个关于\(\partial_t\Psi\)线性的项、一个涉及\(\alpha_i\)的空间导数双线性项，以及质量双线性项\(m\Psi^\dagger\beta\Psi\)。此处只需要这一结构事实。\editorialnote{在尝试的狄拉克拉格朗日量中，符号、括号以及\(i\)的位置在原始材料中明确未定。不从中推断任何依赖于约定俗成的精确公式。}

\begin{figure}[htbp]
\centering
\includegraphics[width=0.76\textwidth]{lecture_09_figure_05.png}
\caption{狄拉克场讨论中的时间导数项、空间导数项和质量项。}
\lectureframe{00:25:37}
\end{figure}

\begin{classroomqa}
\lecturetimestamp{00:26:37}
\audiencequestion{在尝试的狄拉克拉格朗日量中，括号的位置是否正确，即\(i\)是否乘整个表达式？}
\lecturerresponse{此处并不自信地确定其位置，而且对于结构要点而言也不需要：存在一个能生成狄拉克方程的拉格朗日量。}
\end{classroomqa}

现在加入一个标量场\(\phi\)。朴素乘积\(\Psi^\dagger\Psi\,\phi\)不是洛伦兹标量。插入与质量项中出现的同一狄拉克矩阵，得到
\begin{equation}
\mathcal L_{\mathrm{int}}
\propto
\Psi^\dagger\beta\Psi\,\phi.
\end{equation}
等价地，\(\Psi^\dagger\beta\Psi\)是标量狄拉克双线性。对完整拉格朗日量变分，会在狄拉克方程中增加一个正比于\(\phi\)的项，并在标量方程中增加一个正比于\(\Psi^\dagger\beta\Psi\)的源项。这两个场相互散射。

\begin{classroomqa}
\lecturetimestamp{00:28:03}
\audiencequestion{在标量--费米子相互作用中插入的\(\beta\)，是否就是狄拉克质量项中出现的同一个\(\beta\)？}
\lecturerresponse{是的。它是同一个狄拉克矩阵。}
\end{classroomqa}

比较
\begin{equation}
m\Psi^\dagger\beta\Psi
\qquad\hbox{and}\qquad
\Psi^\dagger\beta\Psi\,\phi
\end{equation}
使耦合看起来像是一个被标量场替换的质量参数。当考虑一个具有非零真空值的场时，这一观察将在接近末尾处再次出现。

到目前为止，解释一直是经典的：波包散射其他波包。当这些场成为算符时，同一套代数具有更清晰的粒子含义。

\section{量子场的幂次能做什么}

一个实标量量子场同时包含产生算符和湮灭算符。示意性地，
\begin{equation}
\phi(x)\sim A^+(x)+A^-(x),
\end{equation}
其中\(A^+\)产生一个量子，\(A^-\)湮灭一个量子。因此，\(\phi^3(x)\)包含四种局域过程：
\begin{equation}
\phi^3
\sim
(A^+)^3+(A^+)^2A^-+A^+(A^-)^2+(A^-)^3,
\end{equation}
其中在此示意性解读中，排序细节和数值系数被省略。四种粒子数变化为
\begin{equation}
3\longrightarrow0,
\qquad
0\longrightarrow3,
\qquad
2\longrightarrow1,
\qquad
1\longrightarrow2.
\end{equation}
三个场因子都作用于同一点\(x\)。因此，单个三次项包括三粒子湮灭、三粒子产生、两个量子融合为一个，以及一个量子分裂为两个。在物理振幅中，其中一些局域解读之后可能会被能量--动量守恒移除，但它们在算符乘积本身中全都存在。

二次自由项具有不同的作用。首先抑制相邻点\(x\)与\(x'\)之间差值的归一化：
\begin{equation}
\partial_x\phi
\ \propto\
\phi(x)-\phi(x').
\end{equation}
平方并展开后得到项结构
\begin{align}
\frac12(\partial_x\phi)^2
&\ \propto\
\frac12\bigl[\phi(x)-\phi(x')\bigr]^2
\\
&=
\frac12\phi(x)^2
+\frac12\phi(x')^2
-\phi(x)\phi(x').
\end{align}
每个同点平方都有若干算符解读：它可以在该同一点产生两个量子、湮灭两个量子，或湮灭并重新产生一个量子。这些项不输运粒子。交叉项可以在\(x'\)湮灭一个量子并在\(x\)产生它，或反过来。因此它充当最近邻跃迁项。重复因子将粒子从\(x\)带到\(x'\)，然后到\(x''\)，并继续向前。二次导数项支配自由传播；三次及更高次项散射、产生和湮灭粒子。

\begin{classroomqa}
\lecturetimestamp{00:36:07}
\audiencequestion{导数的有限差分表示是否应包含间隔\(x-x'\)？}
\lecturerresponse{是的。将间隔称为\(L\)，在差值平方后给出因子\(L^{-2}\)。该数值因子改变系数，但不改变跃迁解释。}
\end{classroomqa}

恢复间隔后得到
\begin{align}
\frac12(\partial_x\phi)^2
&\longrightarrow
\frac{1}{2L^2}\bigl[\phi(x)-\phi(x')\bigr]^2
\\
&=
\frac{1}{2L^2}\phi(x)^2
+\frac{1}{2L^2}\phi(x')^2
-\frac{1}{L^2}\phi(x)\phi(x').
\end{align}

质量项也是二次的，但不含导数：
\begin{equation}
\frac12m^2\phi(x)^2.
\end{equation}
其单粒子解读是一个入射量子在同一点被吸收并重新发射。它不移动该量子；它对该量子在此处的存在赋予权重，或在这个意义上计数。同一算符平方也包含其双产生和双湮灭解读。因此，与运动相关的区别在于连接相邻点的交叉项与作用于单点内部的质量插入之间。

\section{量子电动力学的局域顶点}

这些基本粒子过程的草图是费曼图的开端。核心例子是量子电动力学，即关于电子、正电子和光子的理论。

狄拉克方程具有正能模和负能模。用较早的说法，\(\Psi\) 包含两类模式的湮灭算符，而 \(\Psi^\dagger\) 包含相应的产生算符。去掉一个负能电子被重新解释为产生一个正能正电子。经过这种重新标记后，有用的示意性分解为
\begin{align}
\Psi
&\sim
C^-_{e^-}+C^+_{e^+},
\\
\Psi^\dagger
&\sim
C^+_{e^-}+C^-_{e^+}.
\end{align}
因此，\(\Psi\) 湮灭一个电子或产生一个正电子，而 \(\Psi^\dagger\) 产生一个电子或湮灭一个正电子。这种简写不意味着任何归一化、自旋标记或动量积分。

QED 相互作用将狄拉克流与电磁矢量势耦合起来。在本书所用的分量记号中，其结构为
\begin{equation}
\mathcal L_{\mathrm{QED,int}}
\sim
e\left(
\Psi^\dagger\Psi\,A_0
+
\Psi^\dagger\boldsymbol{\alpha}\Psi\mathbin{\boldsymbol{\cdot}}\mathbf A
\right),
\end{equation}
其中
\begin{equation}
J^0=\Psi^\dagger\Psi,
\qquad
\mathbf J=\Psi^\dagger\boldsymbol{\alpha}\Psi.
\end{equation}
符号 \(\sim\) 很重要：时间项与空间项之间的相对符号取决于度规和矢量势约定。不依赖于约定的陈述是，四维流与四维势缩并。\editorialnote{该来源确定了流的空间分量和时间分量，但没有为分量缩并确定度规约定。}

光子场 \(A_\mu\) 本身包含光子的产生算符和湮灭算符。它的矢量性质还携带光子的极化---一个小的方向性 ``标志''，此外还有它的动量。

\begin{classroomqa}
\lecturetimestamp{00:44:28}
\audiencequestion{当 \(A_\mu\) 是矢量，而费米子因子起初看起来是标量时，相互作用怎么可能是标量？}
\lecturerresponse{费米子也提供一个四维矢量：\(\Psi^\dagger\Psi\) 是它的时间分量，\(\Psi^\dagger\boldsymbol\alpha\Psi\) 是它的空间分量。将该流与 \(A_\mu\) 缩并就产生一个标量。}
\end{classroomqa}

将 \(\Psi^\dagger\)、\(\Psi\) 和 \(A_\mu\) 的算符内容相乘，就得到同一个局域顶点的四种常见读法：
\begin{enumerate}
\item 一个电子发射一个光子；
\item 一个电子吸收一个光子；
\item 一个电子和一个正电子湮灭成一个光子；
\item 一个光子产生一个电子--正电子对。
\end{enumerate}
对于第一种读法，\(\Psi\) 湮灭入射电子，\(\Psi^\dagger\) 产生出射电子，而光子场产生光子。选择光子湮灭部分就得到吸收。选择正电子部分就得到对湮灭或对产生。一个局域项包含了整个家族。

在费米子流图中，正电子线上的箭头指向与正电子的传播方向相反。按照这一约定，费米子箭头连续地穿过每一个允许的 QED 顶点。这四种读法的差别在于把外线在入射侧和出射侧之间移动，而不是断开费米子流线。

这些读法中的每一种都带有相同的系数 \(e\)，即所用微观约定中的电荷。这个系数是电荷 \(e\)，而不是精细结构常数，尽管二者密切相关。顶点振幅正比于 \(e\)；在计入所有运动学和归一化因子之后，相应的概率尺度正比于 \(e^2\)。

\begin{classroomqa}
\lecturetimestamp{00:49:44}
\audiencequestion{光子到电子--正电子的读法是否与该顶点的其他读法具有相同的系数？}
\lecturerresponse{是的。它带有相同的电荷系数 \(e\)。}
\end{classroomqa}

\begin{classroomqa}
\lecturetimestamp{00:49:59}
\audiencequestion{既然概率没有单位，是否必须有一个标度因子？}
\lecturerresponse{在所用的自然微观约定中，电荷 \(e\) 是无量纲的。它是振幅尺度，而它的平方给出概率尺度。}
\end{classroomqa}

一个具体的图像是电子撞击墙壁并突然停下。存在某种概率，使能量和动量以光子的形式离开。该辐射振幅中的基本电磁因子是 \(e\)，因此其概率尺度包含 \(e^2\)。这一陈述分离出耦合因子；实际概率还取决于电子的状态和运动学。

\begin{figure}[htbp]
\centering
\includegraphics[width=0.78\textwidth]{lecture_09_figure_06.png}
\caption{\(\Psi\) 和 \(\Psi^\dagger\) 的电子与正电子算符内容，并附有一个电子--光子相互作用草图。}
\lectureframe{00:50:39}
\end{figure}

\begin{classroomqa}
\lecturetimestamp{00:51:15}
\audiencequestion{电荷有单位吗？}
\lecturerresponse{在此处所用的自然微观约定中，没有。宏观电荷常以库仑报告，但粒子理论是相对于选定的基本单位来测量电荷的。在普朗克单位或其他自然单位中，这些量被表示为纯数。}
\end{classroomqa}

这一对比是约定与尺度之别。库仑便于描述导线中的电流；无量纲耦合便于描述微观振幅。\editorialnote{来自00的一个简短问题：52：37--00：53：15关于``分组''带电荷项的表述严重受损，未能引出确定的构造。}

基本QED相互作用中的三个场都在同一个时空点上求值。顶点本身不带有任何内部时间延迟。

\begin{classroomqa}
\lecturetimestamp{00:53:19}
\audiencequestion{电子的吸收与再发射以及光子的发射是同时发生的，还是方程中内置了延迟？}
\lecturerresponse{在基本顶点处，所有算符都在同一个时空点上求值。因此，基本过程中这些作用之间不包含延迟。}
\end{classroomqa}

\begin{classroomqa}
\lecturetimestamp{00:53:57}
\audiencequestion{每个守恒量都必须分别守恒吗？}
\lecturerresponse{是的。电荷就是必须分别守恒的量之一。}
\end{classroomqa}

电荷守恒在图中和代数中都可见。一个把入射电子变成出射正电子的过程会反转电荷，因而不是允许的QED顶点。当时间向上绘制时，电子箭头沿时间向上，正电子箭头向下；有效顶点总是具有一条贯穿的费米子箭头。在代数上，相互作用包含一个\(\Psi\)和一个\(\Psi^\dagger\)。前者在湮灭一个电子或产生一个正电子时使电荷升高一个单位，后者使其降低一个单位。它们的效果相互抵消。

\begin{classroomqa}
\lecturetimestamp{00:56:05}
\audiencequestion{方框中的电子--光子表达式就是整个QED拉格朗日量吗？}
\lecturerresponse{不是。它是相互作用项。支配传播的二次自由光子项和二次狄拉克项尚未写出。}
\end{classroomqa}

\begin{classroomqa}
\lecturetimestamp{00:56:35}
\audiencequestion{经典拉格朗日量是被推导出来的，而量子拉格朗日量是被猜测并检验的，并且动能减势能的结构在量子力学中仍然保留吗？}
\lecturerresponse{动能减势能的模式仍然保留：时间导数提供动能形式，空间梯度将其推广，未微分的场项提供势能形式。量子拉格朗日量是更基本的描述。在包含极大量子的态中，同一个拉格朗日量给出经典波动方程。候选拉格朗日量借助物理原理和更深层的理论被提出，然后通过碰撞实验被确立或排除。}
\end{classroomqa}

大占据数极限解释了为什么同一个表达式既能支配量子粒子理论，又能支配其经典波。这种一致并非事后强加的另一个奇迹；经典场是大量子的集体行为。反过来，实验碰撞足够直接地暴露基本顶点，从而以极高精度确定QED和标准模型的相互作用。

\section{交叉与时空守恒定律}

局部顶点的腿可以在入射位置与出射位置之间交叉。一个被移到入射侧的外出粒子被重新解释为其反粒子，反之亦然。在这个意义上，同一个电子--正电子--光子算符具有发射、吸收、对湮灭、对产生，甚至全入射的局部读法。

\begin{classroomqa}
\lecturetimestamp{00:59:34}
\audiencequestion{如果一个基本顶点图被旋转，它是否仍然是一个有效的局部读法？}
\lecturerresponse{是的。它的外腿可以在入射位置与出射位置之间移动，并伴随相应的粒子--反粒子重新解释。}
\end{classroomqa}

局部算符读法还不是一个完整的物理过程。它的振幅必须对顶点可能发生的每一个时空点进行积分。时间积分强制能量守恒，空间积分强制动量守恒。一个电子、一个正电子和一个光子全部消失进真空的图像作为一种形式上的局部读法是存在的，但积分后的振幅为零，因为该过程无法守恒能量和动量。

\begin{classroomqa}
\lecturetimestamp{01:00:30}
\audiencequestion{全部入射的图像是否意味着一个电子、一个正电子和一个光子可以消失为虚无？}
\lecturerresponse{这是对局域算符的一种解读，但物理振幅包含对顶点位置的积分。这些积分施加能量和动量守恒，因此这个运动学上被禁戒的过程振幅为零。}
\end{classroomqa}

\begin{classroomqa}
\lecturetimestamp{01:01:57}
\audiencequestion{反向的真空到电子--正电子--光子的解读是否存在，被排除的过程在积分后是否会消失？}
\lecturerresponse{反向的局域解读同样存在。对空间和时间的积分会移除任何与守恒定律不相容的过程。}
\end{classroomqa}

拉格朗日量中的导数符合算符展开，而不会对产生和湮灭符号本身求导。场具有如下示意形式
\begin{equation}
\Psi(x)
\sim
\int d^3k\,
\bigl[
C(\mathbf k)e^{i\mathbf k\cdot\mathbf x-i\omega t}
+\cdots
\bigr].
\end{equation}
算符\(C(\mathbf k)\)是这个时空展开中的系数。导数作用于平面波，带下\(\mathbf k\)或\(\omega\)。

\begin{classroomqa}
\lecturetimestamp{01:02:27}
\audiencequestion{对场展开求导时，产生算符本身是否被求导？}
\lecturerresponse{不。在这个展开中，产生和湮灭算符与时空坐标无关。导数作用于指数因子，带下动量和频率。}
\end{classroomqa}

\section{核子与{\proseunicodefont π}介子作为有效理论}

同样的语言可以描述质子、中子和{\proseunicodefont π}介子。质子和中子是费米子，分别由场\(\Psi_p\)和\(\Psi_n\)表示。在这个相互作用计数讨论中，{\proseunicodefont π}介子场被当作三种电荷变体的标量场，
\begin{equation}
\pi^+,
\qquad
\pi^-,
\qquad
\pi^0.
\end{equation}
质子和\(\pi^+\)带一个正单位电荷，\(\pi^-\)带一个负单位电荷，而中子和\(\pi^0\)是电中性的。在更深层次上，{\proseunicodefont π}介子是夸克--反夸克系统，但这种内部描述最初可以不予解决。

一个示意性的、守恒电荷的核子--{\proseunicodefont π}介子场乘积列表是
\begin{equation}
\mathcal L_{N\pi}
\sim
\left\{
\Psi_p^\dagger\Psi_n\,\pi^+,\;
\Psi_n^\dagger\Psi_p\,\pi^-,\;
\Psi_p^\dagger\Psi_p\,\pi^0,\;
\Psi_n^\dagger\Psi_n\,\pi^0
\right\}.
\end{equation}
这里只展示场内容和电荷分配；相对系数、符号以及旋量矩阵结构在此并未固定。第一个{\proseunicodefont π}介子场上的电荷通过检验一个具体解读来确定：一个中子与一个正{\proseunicodefont π}介子一起可以产生一个质子，而在那里使用负{\proseunicodefont π}介子则会把负的入射电荷变成正出射电荷。共轭项给出相反的带电过程。\editorialnote{{\proseunicodefont π}介子的符号在源中经过多次修正，才达到这里所示的守恒电荷配对；所展示的顺序记录了那个最终的示意结构。}

每一项同样有若干交叉解读。把一个出射中子移到入射侧会把它变成一个入射反中子。因此带电顶点可以描述一个质子和一个反中子结合成一个\(\pi^+\)，以及在运动学允许时诸如\(p\to n+\pi^+\)的过程。中性项让质子或中子发射或吸收一个\(\pi^0\)。电荷守恒选出允许的家族，而能量、动量以及任何进一步的对称性决定哪一个成员能在特定实验中发生。

在夸克子结构为人所知之前，这是一个有用的有效描述。每当探针无法分辨该结构时，它仍然有用。质子具有有限的组分分布，有点像原子具有内部结构，尽管尺度小得多。如果实验转移的动量太小，不足以打碎或分辨质子，质子可以被当作具有自身场的点粒子。足够硬的碰撞会暴露夸克，此时点粒子质子场就不再充分。

因为质子带电，它还具有局域的质子--光子相互作用，示意为\(J_p^\mu A_\mu\)，其中\(J_p^\mu\)是质子流。电子有相应的电子--光子相互作用。质子不能改为通过发射一个光子而变成一个中子：\(p\to n+\gamma\)违反电荷守恒。

\begin{classroomqa}
\lecturetimestamp{01:10:28}
\audiencequestion{当一个质子和一个电子相互作用时会发生什么？}
\lecturerresponse{每个带电粒子都有一个光子顶点。质子可以发射一个光子而仍保持为质子，电子可以吸收该光子而仍保持为电子。质子--电子散射由这两个局域顶点以及光子在其间的传播组装而成。}
\end{classroomqa}

\section{局域性与重复的基本过程}

对于完整的质子--电子交换，并不存在特殊的非定域项。一个定域顶点在质子处发射光子，二次光子项将其从一个邻近点带到另一个邻近点，另一个定域顶点在电子处吸收它。这一扩展过程是那些基本片段的乘积。
\begin{equation}
\begin{aligned}
\bigl[J_p^\mu A_\mu\bigr](x_p)
&\longrightarrow
\underbrace{x_p\to x_1\to x_2\to\cdots\to x_e}_{\text{quadratic photon propagation}}
\\
&\longrightarrow
\bigl[J_e^\nu A_\nu\bigr](x_e).
\end{aligned}
\end{equation}
这个示意链中的箭头表示相继的邻近点插入，而不是一种新的非定域相互作用。

\begin{classroomqa}
\lecturetimestamp{01:12:06}
\audiencequestion{人们会如何写出整个质子--电子交换的拉格朗日量？}
\lecturerresponse{人们不会为整个扩展事件添加一项。它由一个定域的质子--光子顶点、重复的二次传播项以及一个定域的电子--光子顶点构成。}
\end{classroomqa}

定域结构有两种形式。未加区分的相互作用项作用于一个时空点。以有限差分表示的导数项只将该点与一个邻近点连接起来。这两种形式都不会把粒子直接转移过一段有限距离。长程传播由许多无穷小移动组装而成；不存在基本的作用于距离。

\begin{classroomqa}
\lecturetimestamp{01:14:36}
\audiencequestion{拉格朗日量只包含一个点上的项吗？}
\lecturerresponse{相互作用项作用于一个点。导数项将一个点与一个无穷小邻近点连接起来。发射、许多邻近跳跃和吸收共同产生一次有限交换。}
\end{classroomqa}

\begin{classroomqa}
\lecturetimestamp{01:15:36}
\audiencequestion{扩展过程的所有片段都存在吗？把它们复合起来是否意味着相乘？}
\lecturerresponse{是的。每一次基本发射、跳跃或吸收都编码在拉格朗日量中。重复插入的乘积将这些片段复合成更大的过程。}
\end{classroomqa}

所讨论的相对论量子场论预设了其场被定义于其上的时空背景。不应将其与关于动态时空几何的完整量子理论相混淆。

\begin{classroomqa}
\lecturetimestamp{01:16:28}
\audiencequestion{这种复合在强相对论情形下，例如在黑洞视界附近，是否仍然有效？}
\lecturerresponse{黑洞的推广超出了这一构造的范围。这里所使用的量子场论本身并不提供对黑洞动力学的完整描述，而在那里引力是必不可少的。}
\end{classroomqa}

场是基本定域作用的算符。一个场因子可以吸收一个粒子或发射一个粒子。两个因子的乘积可以在同一点吸收一个粒子并发射另一个粒子；更高阶的乘积可以进一步改变粒子内容。定域算符的乘积将这些作用连接成历史。在时空中，运动被表示为在一个小区域中的吸收、向邻近区域的再发射、再一次吸收和再一次再发射。让步长趋于零而步数增加，就给出连续轨迹。

\begin{classroomqa}
\lecturetimestamp{01:17:30}
\audiencequestion{在这种定域步骤图像中，场、距离和时间应如何理解？}
\lecturerresponse{场是定域吸收和发射的算符。乘积将这些基本作用连接起来，而穿越时空的运动由无穷小邻近步骤构成，然后才取连续极限。}
\end{classroomqa}

因此，一个正规化的起点是把时空划分为小单元，场算符附着在这些单元上。\editorialnote{在01处的听众措辞：17：30--01：20：40部分被遮蔽；转述的问题只保留了其可恢复的主题，而关于定域算符和相继时空步骤的回答是清楚的。}导数变成差分。二次交叉项连接相邻单元，而质量和相互作用项作用于一个单元内部。这些单元是定义极限理论的工具，而不是额外的粒子或新的物理质量标度。

狭义相对论效应不必在每一步单独添加。当拉格朗日量是洛伦兹标量时，由它组装出的振幅是洛伦兹不变的，而诸如时间膨胀之类的效应就由这一结构得出。

\begin{classroomqa}
\lecturetimestamp{01:21:10}
\audiencequestion{高能运动和时间膨胀会在多步过程中造成额外问题吗？}
\lecturerresponse{不需要单独的修正。如果拉格朗日量是洛伦兹标量，诸如时间膨胀之类的狭义相对论效应已经内建于结果之中。除此之外，该问题背后的确切关切并不完全清楚。}
\end{classroomqa}

产生和湮灭算符既可以在位置基中表示，也可以在动量基中表示。位置空间的产生算符是动量空间产生算符的傅里叶叠加，反之亦然。不存在一个算符能产生同时具有精确指定的位置和精确指定的动量的态。将粒子高度局域化需要宽范围的动量；将其在动量上高度局域化则会使其在位置上弥散开来。两个基之间的傅里叶关系正是这一语言下的不确定性原理。

\begin{classroomqa}
\lecturetimestamp{01:22:48}
\audiencequestion{产生和湮灭算符与位置、动量以及不确定性原理有何关系？}
\lecturerresponse{它们既可以表示为位置的函数，也可以表示为动量的函数。这两种表示互为傅里叶变换。一个高度局域化的位置态需要许多动量，而一个高度局域化的动量态则在位置上弥散开来；不存在能同时产生精确的\(x\)和\(p\)的产生算符。}
\end{classroomqa}

现在可以完整地陈述元胞构造了。用\(n\)标记各元胞，将导数替换为有限差分，例如
\begin{equation}
\partial_\mu\phi(x)
\longrightarrow
\frac{\phi_{n+\hat\mu}-\phi_n}{a},
\end{equation}
并将每一种场都放入每一个元胞中。一个二次交叉项从一个元胞中移除一个量子，并在相邻元胞中产生一个量子。在受调控的表达式被定义之后，令间距\(a\)趋于零，就得到连续理论。
\begin{equation}
\underbrace{\phi_n\phi_{n+\hat\mu}}_{\text{hop }n\leftrightarrow n+\hat\mu},
\qquad
\underbrace{m^2\phi_n^2}_{\text{same-cell mass insertion}}.
\end{equation}

\begin{classroomqa}
\lecturetimestamp{01:25:16}
\audiencequestion{将时空细分为元胞的构造能够完成吗？}
\lecturerresponse{能。场以元胞为指标，导数变为有限差分，二次项使粒子在相邻元胞之间跳跃。当元胞尺寸趋于零时，就恢复了连续情形。}
\end{classroomqa}

电子通过其二次项在元胞之间跳跃，光子也是如此。当一个电子和一个光子在同一个元胞中相遇时，QED相互作用可以吸收并重新发射该电子，同时产生或湮灭光子。标量质量项有一种同元胞解读：一个粒子进入并离开该元胞，其系数正比于\(m^2\)。它并不提供跳跃。

\begin{classroomqa}
\lecturetimestamp{01:28:18}
\audiencequestion{元胞尺寸与质量系数有关吗？}
\lecturerresponse{无关。元胞尺寸是一种最终要取为零的调控间距；质量系数则是拉格朗日量中的一个独立参数。}
\end{classroomqa}

\begin{classroomqa}
\lecturetimestamp{01:28:43}
\audiencequestion{将局域因子相乘，是一种在时空上的联合振幅构造，而不仅仅是时间上的一个序列吗？}
\lecturerresponse{是的。各元胞覆盖时空，它们的局域贡献在整个时空网格上组合起来。该乘积为所需的局域事件和传播步骤提供了联合振幅。}
\end{classroomqa}

\section{源、探测器以及它们之间的动力学}

一个完整过程是时空中的一段历史。它的初始条件由实验者所制备的东西确定，末态条件则由探测器所记录的东西确定。一个简单的源是发射电子的热灯丝。灯丝处的一个场算符产生该入射电子；探测器位置处的场算符描述它最终的吸收。

\begin{classroomqa}
\lecturetimestamp{01:29:30}
\audiencequestion{这样一个时空过程的初始阶段是什么？}
\lecturerresponse{它是在源处制备的。由热灯丝发射的电子是一个简单例子：灯丝处的一个场算符产生该电子，而探测器处的场算符则指定最终的吸收事件。}
\end{classroomqa}

同样的构造也可以从真空出发来陈述。源处的产生算符在真空上构建入射态。与探测器相关的湮灭算符选出出射态。这些算符指定边界条件；它们并不取代源与探测器之间整个区域内的动力学。

\begin{classroomqa}
\lecturetimestamp{01:31:11}
\audiencequestion{把动力学只放在发射体处，而不是放在整个系统中，会是错误的吗？}
\lecturerresponse{是的。从真空出发，用源算符制备入射粒子，用探测器算符选出最终粒子。在它们之间，反复插入拉格朗日量，在整个时空上起作用，移动粒子、产生粒子并湮灭粒子。}
\end{classroomqa}

有了足够多的插入，若干粒子可以同时穿过许多相邻元胞；粒子也可以沿途被产生或湮灭。初始算符说明什么进入，最终算符说明什么被记录，而反复的拉格朗日量操作则提供它们之间的一切。

\begin{classroomqa}
\lecturetimestamp{01:33:24}
\audiencequestion{如果任意多个拉格朗日插入都可以出现，那么它们的位置是否被积分掉？}
\lecturerresponse{是的。它们的时空位置被积分掉。这里并不推导精确的费曼规则；眼下要点在于，简单的局域场表达式编码了原初相互作用，而完整过程正是由这些相互作用组装而成的。}
\end{classroomqa}

\section{相位对称性、电荷与振幅}

把原初过程写进拉格朗日量，也使它们的对称性变得可见。考虑电子场的整体相位变换，
\begin{equation}
\Psi\longrightarrow e^{i\theta}\Psi,
\qquad
\Psi^\dagger\longrightarrow e^{-i\theta}\Psi^\dagger,
\qquad
A_\mu\longrightarrow A_\mu.
\end{equation}
含有等量\(\Psi\)与\(\Psi^\dagger\)的项保持不变，因为相位相互抵消。数目不等的项则会获得一个相位。这些数目相等，正是电荷守恒的代数形式。

这是拉格朗日量描述的一个核心优势。它的场因子展示出基本过程，它们的组合展示出对称性，而对称性又展示出守恒定律。反复插入把制备好的初态演化到选定的末态。结果是一个量子振幅。精确计算需要费曼规则，但其结构内容已经可见：局域操作在守恒定律所允许的每一阶中反复出现。

\section{场的测量与对易子}

场算符保留算符通常的量子力学含义。它们的本征值是场测量的可能结果，而它们的本征矢描述与这些值相联系的态。要探测电场，可以把一个带电粒子放入其中，并观察该粒子如何加速。

\begin{classroomqa}
\lecturetimestamp{01:37:06}
\audiencequestion{场算符的本征值和本征矢是否发挥它们通常的量子力学作用？}
\lecturerresponse{是的。本征值是场可能被测量到的值，而场可以通过它们对粒子的效应来探测。对易子决定同时测量的限制。}
\end{classroomqa}

平凡恒等式
\begin{equation}
[\phi(x),\phi(x)]=0
\end{equation}
并不能解决测量问题。必须指定两个可观测量以及它们的时空关系。一个场与其共轭动量构成一对正则非对易量，而类时相关点上的场算符不必对易；类空间隔上的局域玻色场则确实对易，正如定域性所要求的那样。因此，对易子与不确定性之间的熟悉联系仍然成立，但只有在正确识别出算符对之后才成立。\editorialnote{原文对异点处非对易性的陈述过于宽泛。此处的限定保留了其测量要点，同时强制实现类空间隔局域玻色场的微观因果性。}

\begin{classroomqa}
\lecturetimestamp{01:39:57}
\audiencequestion{电场或电磁场是以伏特测量，还是以伏特每米测量？}
\lecturerresponse{电场以伏特每米测量。伏特和米并不是描述微观场过程最有用的语言，而微观场过程更自然地用光子来组织。反过来说，量子场论也不是决定如何给房子浇水的有用语言。}
\end{classroomqa}

对一种描述自然的单位，未必是对另一种描述有启发性的单位。伏特和伏特每米适合宏观仪器；光子的产生和湮灭适合微观过程。

\begin{classroomqa}
\lecturetimestamp{01:41:10}
\audiencequestion{核心思想是否是：拉格朗日量中的项是组合成更大过程的构建块？}
\lecturerresponse{是的。基本过程从一个点到下一个点反复出现。每个基本要素作用于一个时空点，或通过导数项连接相邻点。表面上非局域的传播正是由这些局域步骤构建起来的；不存在基本的作用于远距离。}
\end{classroomqa}

\section{局域过程的乘积与路径求和}

有一种紧凑的启发式方式可以展示局域过程的重复。把时空分成胞元，并写出
\begin{equation}
\prod_x\bigl(1+\mathcal L_I(x)\bigr),
\end{equation}
其中\(x\)标记胞元，\(\mathcal L_I(x)\)是局域相互作用拉格朗日量。这并不是精确的连续统公式。它是一个组合式的替代表示，揭示出乘积如何生成局域事件的每一种允许选择。

\begin{classroomqa}
\lecturetimestamp{01:42:18}
\audiencequestion{许多产生、传播和湮灭步骤能否压缩进一个表达式？}
\lecturerresponse{从启发式上说，对每个时空胞元乘上一个因子\(1+\mathcal L_I(x)\)。展开该乘积会生成没有事件、在任一选定点有一个事件，或在若干点有事件的情形。}
\end{classroomqa}

对于两个格点，
\begin{align}
\bigl(1+\mathcal L_I(x)\bigr)
\bigl(1+\mathcal L_I(x')\bigr)
&=
1+\mathcal L_I(x)+\mathcal L_I(x')
\notag\\
&\quad{}+\mathcal L_I(x)\mathcal L_I(x').
\end{align}
第一项表示两个相互作用都不发生。接下来的两项分别在\(x\)处或\(x'\)处放置一个基本相互作用。最后一项在两个点处都放置相互作用。当因子很多时，会出现涉及一个点、两个点、三个点等等的项。一个含有多个顶点的项可以描述例如两个入射粒子产生三个出射粒子。

二次传播片段也进入同样的构造。有限位移由一系列相邻插入的链条表示。在展开的乘积中，某处会出现位于一个格点、其相邻格点、再下一个格点等等的因子，将粒子从一个端点携带到另一个端点。

\begin{classroomqa}
\lecturetimestamp{01:46:29}
\audiencequestion{如果有限传播需要许多局域运动项，那么所需的运动是如何恢复的？}
\lecturerresponse{这个大的乘积包含一系列相邻插入的链条，将粒子从其初始点携带到其最终点。通常存在许多沿不同路径的这种链条。每条路径贡献一个振幅，总振幅是它们的和。}
\end{classroomqa}

如果\(r\)标记一条路径，那么示意性地
\begin{equation}
\mathcal A_{\mathrm{total}}=\sum_r\mathcal A_r,
\qquad
P=\left|\mathcal A_{\mathrm{total}}\right|^2.
\end{equation}
并不是先给路径分配普通概率再相加。而是先相加振幅：
\begin{equation}
\left|\sum_r\mathcal A_r\right|^2
=
\sum_r|\mathcal A_r|^2
+
\sum_{r\ne s}\mathcal A_r\mathcal A_s^*.
\end{equation}
交叉项就是干涉。

\begin{classroomqa}
\lecturetimestamp{01:48:27}
\audiencequestion{这与双缝实验一致吗？}
\lecturerresponse{完全一致。通过两条缝的备选路径贡献振幅。对其和取平方会产生干涉项，正如局域步构造中许多时空路径相互干涉一样。}
\end{classroomqa}

\begin{classroomqa}
\lecturetimestamp{01:49:00}
\audiencequestion{每个因子\(1+\mathcal L_I\)中的\(1\)从何而来？}
\lecturerresponse{它是一个组合装置，使一个格点既可以贡献一个相互作用，也可以不贡献相互作用。它与路径积分方法有关，但所展示的乘积只是对技术性表达式的一种启发式近似。}
\end{classroomqa}

精确的形式体系提供了恰当的时空积分、排序、归一化和相位。这个乘积之所以有用，是因为它揭示了各种组合，而不假装自己就是完整的形式体系。

\begin{classroomqa}
\lecturetimestamp{01:49:57}
\audiencequestion{这是稳定作用量原理吗？}
\lecturerresponse{不是。经典稳定作用量选出一条特殊的经典轨迹。量子力学则是对轨迹求和振幅，其权重由作用量决定。即使是一条极其间接的路径---从办公室到附近一所房子的途中经过巴黎---也属于量子求和。}
\end{classroomqa}

对比很简单，尽管关系很微妙。经典力学选出一条稳定路径。量子力学对路径求和；在适当的极限下，相位在经典轨迹附近相互增强。

\section{推断拉格朗日量与选择部门}

一旦给定拉格朗日量和指定的过程，计算其概率的规则就是确定的。计算可能仍然需要巧思，正如困难的数学往往如此，但更深层的不确定性在于发现自然界使用哪些场和耦合。

\begin{classroomqa}
\lecturetimestamp{01:50:43}
\audiencequestion{艺术在于表达拉格朗日量，还是在于使用它？}
\lecturerresponse{主要的艺术在于从实验数据推断拉格朗日量。一旦它已知，诸如两个粒子进入、四个粒子离开这样的过程就有一个确定的振幅，可以用精确规则计算。执行数学运算可能需要聪明才智，但这不是理论中的歧义。}
\end{classroomqa}

粒子碰撞暴露出基本顶点。人们提出各项，往往在对称性和更深层原理的指导下，并将得到的振幅与观测结果进行检验。成功既确定了场的允许组合，也确定了它们的数值系数。

\begin{classroomqa}
\lecturetimestamp{01:52:24}
\audiencequestion{每一种相互作用都有一个单独的拉格朗日量吗---一本大的拉格朗日量之书？}
\lecturerresponse{实际图景反而是包含所有场和相互作用的一个大拉格朗日量，尽管这个图景最终可能是错的。可识别的部门作为和中的片段出现。一个新粒子增加一个场；一个新基本过程增加一个相互作用项。}
\end{classroomqa}

完整的表达式可以显得杂乱，却不必长达一本书。作为一个刻意粗略的尺度，人们可以设想大约一百个被称为基本粒子的对象以及数百个项。量子电动力学、强相互作用与弱相互作用，再加上一个引力部分，仍然可能只占几页写得密密麻麻的篇幅。确切的数目并不是重点。实际使用的图景是一个由许多可辨认的部件组装而成的动力学表达式，而且它明确地只是一个最终可能被证明并不完备的图景。

这些场是否决定了一种更为基本的时空几何，超出了此处所使用的标准量子场框架。\editorialnote{关于01:55:00--01:55:20的问题损坏过甚，无法可靠地辨认出流形或几何构造；只有这一范围界限式的回应是可靠的。}

\begin{classroomqa}
\lecturetimestamp{01:55:22}
\audiencequestion{实验测定的那些数就是耦合常数吗？}
\lecturerresponse{耦合常数是拉格朗日量中各项的系数。二次质量项的系数\(m^2\)是一个例子；\(\phi^4\)、\(\phi^6\)或更高次幂的系数则是其他例子。}
\end{classroomqa}

一个系数控制着一族相关的振幅，因为每个场乘积都有若干种产生与湮灭的读法。因此，知道了拉格朗日量所决定的远不止一个孤立的反应。

\begin{classroomqa}
\lecturetimestamp{01:56:05}
\audiencequestion{为了计算量子电动力学，必须知道整个拉格朗日量吗？}
\lecturerresponse{无限高的精度将要求知道电子与光子能够通过其发生相互作用的每一个部分。然而在原子能标下，电子--光子部分通常可以被孤立出来。能量不足以分辨微小的夸克结构，因此其他部分可以在所需精度内被忽略。}
\end{classroomqa}

实际计算使用的是与实验能标相关的那个部分。这类似于流体力学：一个复杂的微观系统可以容许一种精确的粗粒化描述，而不必展示其所有更深层的组分。因此，量子场论可以同时是唯象的又是精确的---在从实验推断出的场与系数方面是唯象的，在这些要素被指定之后所应用的规则方面是精确的。

\begin{classroomqa}
\lecturetimestamp{01:58:09}
\audiencequestion{标准模型包含引力吗，这些相对论性拉格朗日量包含广义相对论吗？}
\lecturerresponse{正式称为标准模型的理论不包含引力，尽管在讨论更广泛的理论时可以把引力附加进去。此处所考虑的场论把量子力学与狭义相对论结合起来。它们本身并不包含广义相对论。}
\end{classroomqa}

\section{作为质量系数的真空值}

标量场的质量表现为\(\phi^2\)的系数。因此，当另一个场具有非零真空值时，与那个场的相互作用就可以模仿一个质量项。

\begin{classroomqa}
\lecturetimestamp{01:59:16}
\audiencequestion{如果粒子质量通过希格斯相互作用产生，那么希格斯耦合能起质量项的作用吗？}
\lecturerresponse{能。项\(H\phi^2\)是三次的，通常表示一个带有两条\(\phi\)腿和一条\(H\)腿的相互作用。如果势\(V(H)\)选出一个\(H\)非零的真空，那个真空值就是一个乘在\(\phi^2\)上的数，因此它贡献一个有效的二次质量项。}
\end{classroomqa}

示意地考虑
\begin{equation}
\mathcal L\supset H\phi^2-V(H).
\end{equation}
作为一个算符项，\(H\phi^2\)可以吸收一个\(\phi\)、发射一个\(\phi\)，并发射或吸收一个\(H\)量子。假设势在\(H=0\)之外有两个极小值，而真空选出其中一个：
\begin{equation}
\langle H\rangle=v\ne0.
\end{equation}
把场写成其真空值加上一个涨落，
\begin{equation}
H=v+h.
\end{equation}
于是
\begin{equation}
H\phi^2=v\phi^2+h\phi^2.
\end{equation}
第一项在\(\phi\)中是二次的，并对其质量系数有贡献；第二项仍然是一个与涨落场\(h\)的相互作用。\editorialnote{来源在陈述非零真空值成为质量系数时结束了这句话。所展示的替换是直接的代数补全；其归一化与符号取决于完整的拉格朗日量约定。}

\begin{classroomqa}
\lecturetimestamp{02:01:29}
\audiencequestion{标准模型包含虚粒子的产生与湮灭吗？}
\lecturerresponse{绝对包含。虚产生与虚湮灭是标准模型拉格朗日量所编码的反复量子过程的内在组成部分。}
\end{classroomqa}
